\documentclass[10pt,a4paper]{amsart}
\usepackage[margin=19mm]{geometry}
\usepackage{amsmath,amssymb,mathtools}
\usepackage[scr=rsfs,cal=euler]{mathalfa}
\usepackage{xfrac}
\usepackage[unicode,hidelinks,bookmarks=false]{hyperref}
\newcommand{\ie}{i.e.\ }
\newcommand{\cf}{cf.\ }
\newcommand{\resp}{respectively\ }
\newcommand{\Subsection}{\subsection}
\providecommand{\nobreakdash}{\nobreak\hbox{-}}
\setlength{\emergencystretch}{2em}
\multlinegap0pt
\usepackage{bm}
\usepackage{bbm}
\usepackage{dsfont}
\usepackage{xspace}
\usepackage{cancel}
\usepackage{mathtools}
\usepackage{amsthm}
\makeatletter
\@ifundefined{theorem}{\newtheorem{theorem}{Theorem}}{}
\makeatother
\newcommand{\N}{\mathbb{N}} %% Conjunto naturales:     \N
\newcommand{\Z}{\mathbb{Z}} %% Conjunto enteros:       \Z
\newcommand{\cB}{\mathcal{B}}
\newcommand{\cO}{\mathcal{O}}
\newcommand{\fS}{\mathfrak{S}}
\title[On the Dynamics of Weighted Composition Operators II]{On the Dynamics of Weighted Composition Operators II}
\author{Nilson C. Bernardes Jr., Antonio Bonilla, Jo\~ao V. A. Pinto}
\date{Source: arXiv:2512.06425v2 (CC BY 4.0)}
\begin{document}
\maketitle

\noindent\emph{Source and licence.} On the Dynamics of Weighted Composition Operators II. Version arXiv:2512.06425v2 (CC BY 4.0), distributed under the Creative Commons Attribution 4.0 International licence, as recorded in the arXiv licence metadata for this exact version and shown on the abstract page https://arxiv.org/abs/2512.06425v2. Full rights notice: licenses/arxiv-2512.06425-CC-BY-4.0.txt.\par\smallskip

\noindent\textit{Editorial scope.} Counted unit: the Theorem whose source label is \texttt{ExpCX-T1} in the article (source0.tex lines 328--348, source label \texttt{ExpCX-T1}), delivered together with the article's own complete original proof of it (source0.tex lines 350--470, one \texttt{proof} environment of 121 lines). No mathematics is written, repaired or re-derived: statement and proof are the article's own text line for line. Required context retained uncounted: none. Every definition, display and label of those lines is retained unchanged. Omitted from this package and not redistributed: 1--327, 349--349, 471--1488. Nothing inside a retained excerpt is omitted or altered: every retained body is delivered exactly as the article's lines are written, so no omission receipt of any kind accompanies this package. Citations: the retained excerpts carry no citation command at all, so no bibliography entry of the article is transcribed and no reference is redistributed. Reference adaptation: none was needed and none was made; every reference inside the delivered unit resolves inside it, so no unresolved reference or citation is produced. Printed slips kept literal: no printed word, symbol, hypothesis or number of the counted statement or of its proof was altered. Layout and class: the article's own class is \texttt{article} with packages including inputenc, geometry, babel, enumitem, amsmath, amsfonts, amssymb, amsthm, amscd, tikz, xcolor, mathrsfs, xy; the wrapper is the lane's pinned \texttt{amsart} editorial wrapper at 10pt on A4, which restates the theorem-like environments the retained text needs and adds the article's own macro definitions that the retained text needs (\texttt{\textbackslash N}, \texttt{\textbackslash Z}, \texttt{\textbackslash cB}, \texttt{\textbackslash cO}, \texttt{\textbackslash fS}), with extra wrapper packages (bm, bbm, dsfont, xspace, cancel, mathtools). The delivered numbering is the unit's own. Assets: no retained span contains a figure environment, a graphics-inclusion command, a PDF-page inclusion, a table or a TikZ picture; no image file is copied into this package, the package declares no asset dependency and no image file is redistributed with it. Licence: version arXiv:2512.06425v2 of this article is distributed under the Creative Commons Attribution 4.0 International licence, as recorded in the arXiv licence metadata for this exact version and shown on the abstract page https://arxiv.org/abs/2512.06425v2. A full rights notice is delivered at licenses/arxiv-2512.06425-CC-BY-4.0.txt. Characters: every retained excerpt is pure ASCII, so the delivered engine, XeLaTeX with its default Latin Modern text font, reports no missing character.
\begin{theorem}\label{ExpCX-T1}
Let $M$ be an admissible subspace of $C_\fS(X;Y)$, $\cB$ a base of the bornology $\fS$ and $\cO$ a base of the topology of $X$,
with $\varnothing \not\in \cB$ and $\varnothing \not\in \cO$. 
For the invertible weighted composition operator $C_{w,f}$ on $M$, the following statements hold:
\begin{itemize}
\item [\rm (a)] $C_{w,f}$ is expansive if and only if, for each $O \in \cO$, there exists $B \in \cB$ such that
    \[
    \sup_{n \in \Z} \|w^{(n)}\|_{B \cap f^{-n}(O)} = \infty.
    \]
\item [\rm (b)] $C_{w,f}$ is average expansive if and only if, for each $O \in \cO$, there exists $B \in \cB$ such that
    \[
    \sup_{n \in \N} \Big(\frac{1}{2n+1} \sum_{j=-n}^n \|w^{(j)}\|_{B \cap f^{-j}(O)}\Big) = \infty.
    \]
\item [\rm (c)] $C_{w,f}$ is uniformly expansive if and only if, for each $A \in \cB$,
    there exists $B \in \cB$ such that we can write $\cO_A = \cO_A^+ \cup \cO_A^-$, where
    \[
    \lim_{n \to \infty} \inf_{O \in \cO_A^+} \|w^{(n)}\|_{B\cap f^{-n}(O)} = \infty \ \ \text{ and } \ \ 
    \lim_{n \to \infty} \inf_{O \in \cO_A^-} \|w^{(-n)}\|_{B\cap f^{n}(O)} = \infty.
    \]
\end{itemize}
\end{theorem}

\begin{proof}
(c): Suppose that $C_{w,f}$ is uniformly expansive. 
Given $(A,\alpha) \in \cB \times I$, there exists $(B,\beta) \in \cB \times I$ such that the unit sphere $S_{\|\cdot\|_{A,\alpha}}$
of the seminorm $\|\cdot\|_{A,\alpha}$ in the space $M$ can be written as
$S_{\|\cdot\|_{A,\alpha}} = S^{+}_{A,\alpha} \cup S^{-}_{A,\alpha}$, where
\begin{equation}\label{ExpCX-Eq1}
\lim_{n \to \infty} \inf_{\varphi \in S^{+}_{A,\alpha}} \|(C_{w,f})^n(\varphi)\|_{B,\beta} = \infty \ \ \text{ and } \ \ 
\lim_{n \to \infty} \inf_{\varphi \in S^{-}_{A,\alpha}} \|(C_{w,f})^{-n}(\varphi)\|_{B,\beta} = \infty.
\end{equation} 
For each $O \in \cO_A$, take a point $x_O \in O \cap A$ and a vector $y_O \in Y$ with $\|y_O\|_\alpha = 1$.
By (C2), there is a continuous function $\phi_O : X \to [0,1]$ such that 
$\phi_O(x_O) = 1$, $\phi_O(X \backslash O) = \{0\}$ and $\varphi_O = \phi_O \otimes y_O \in M$.
Define
\[
\cO_A^+ = \{O \in \cO_A : \varphi_O \in S^{+}_{A,\alpha}\} \ \ \text{ and } \ \
\cO_A^- = \{O \in \cO_A : \varphi_O \in S^{-}_{A,\alpha}\}.
\]
Since $\varphi_O \in S_{\|\cdot\|_{A,\alpha}}$ for all $O \in \cO_A$, we have that
\[
\cO_A = \cO_A^+ \cup \cO_A^-.
\]
For each $O \in \cO_A$ and $n \in \Z$, we have that
\[
\|(C_{w,f})^n(\varphi_O)\|_{B,\beta} = \|w^{(n)} \cdot (\varphi_O \circ f^n)\|_{B \cap f^{-n}(O),\beta} 
  \leq \|w^{(n)}\|_{B \cap f^{-n}(O)} \|y_O\|_\beta.
\]
Hence, for each $n \in \N$,
\begin{equation}\label{ExpCX-Eq2}
\inf_{O \in \cO_A^+} \|w^{(n)}\|_{B\cap f^{-n}(O)} 
  \geq \frac{1}{\|y_O\|_\beta} \inf_{\varphi \in S^{+}_{A,\alpha}} \|(C_{w,f})^n(\varphi)\|_{B,\beta}
\end{equation}
and
\begin{equation}\label{ExpCX-Eq3}
\inf_{O \in \cO_A^-} \|w^{(-n)}\|_{B\cap f^{n}(O)} 
  \geq \frac{1}{\|y_O\|_\beta} \inf_{\varphi \in S^{-}_{A,\alpha}} \|(C_{w,f})^{-n}(\varphi)\|_{B,\beta}.
\end{equation}
By (\ref{ExpCX-Eq1}), (\ref{ExpCX-Eq2}) and (\ref{ExpCX-Eq3}), the limits in (c) are infinite.

Conversely, suppose that the condition in (c) holds. Take $(A,\alpha) \in \cB \times I$.
There exists $B \in \cB$ such that we can write $\cO_A = \cO_A^+ \cup \cO_A^-$, where
\begin{equation}\label{ExpCX-Eq4}
\lim_{n \to \infty} \inf_{O \in \cO_A^+} \|w^{(n)}\|_{B\cap f^{-n}(O)} = \infty \ \ \text{ and } \ \ 
\lim_{n \to \infty} \inf_{O \in \cO_A^-} \|w^{(-n)}\|_{B\cap f^{n}(O)} = \infty.
\end{equation}
For each $\varphi \in S_{\|\cdot\|_{A,\alpha}}$, take $O_\varphi \in \cO_A$ such that
\[
\|\varphi(x)\|_\alpha \geq \frac{1}{2} \ \ \text{ for all } x \in O_\varphi.
\]
Let
\[
S^{+}_{A,\alpha} =\{\varphi \in S_{\|\cdot\|_{A,\alpha}} : O_\varphi \in \cO_A^+\} \ \ \text{ and } \ \
S^{-}_{A,\alpha} =\{\varphi \in S_{\|\cdot\|_{A,\alpha}} : O_\varphi \in \cO_A^-\}.
\]
Clearly, 
\[
S_{\|\cdot\|_{A,\alpha}} = S^{+}_{A,\alpha} \cup S^{-}_{A,\alpha}.
\] 
For each $\varphi \in S_{\|\cdot\|_{A,\alpha}}$ and $n \in \Z$, we have that
\[
\|(C_{w,f})^n(\varphi)\|_{B,\alpha} \geq \|w^{(n)} \cdot (\varphi \circ f^n)\|_{B \cap f^{-n}(O_\varphi),\alpha} 
  \geq \frac{1}{2}\, \|w^{(n)}\|_{B \cap f^{-n}(O_\varphi)}.
\]
Hence, for each $n \in \N$,
\begin{equation}\label{ExpCX-Eq5}
\inf_{\varphi \in S^{+}_{A,\alpha}} \|(C_{w,f})^n(\varphi)\|_{B,\alpha} 
  \geq \frac{1}{2} \inf_{O \in \cO_A^+} \|w^{(n)}\|_{B\cap f^{-n}(O)}
\end{equation}
and
\begin{equation}\label{ExpCX-Eq6}
\inf_{\varphi \in S^{-}_{A,\alpha}} \|(C_{w,f})^{-n}(\varphi)\|_{B,\alpha} 
  \geq \frac{1}{2} \inf_{O \in \cO_A^-} \|w^{(-n)}\|_{B\cap f^n(O)}.
\end{equation}
By (\ref{ExpCX-Eq4}), (\ref{ExpCX-Eq5}) and (\ref{ExpCX-Eq6}),
\[
\lim_{n \to \infty} \inf_{\varphi \in S^{+}_{A,\alpha}} \|(C_{w,f})^n(\varphi)\|_{B,\alpha} = \infty \ \text{ and } \ 
\lim_{n \to \infty} \inf_{\varphi \in S^{-}_{A,\alpha}} \|(C_{w,f})^{-n}(\varphi)\|_{B,\alpha} = \infty,
\]
which proves that $C_{w,f}$ is uniformly expansive.

\smallskip\noindent
(b): Suppose that $C_{w,f}$ is average expansive. Given $O \in \cO$, choose a point $a \in O$ and a vector $y \in Y \backslash \{0\}$.
By (C2), there is a continuous function $\phi : X \to [0,1]$ such that 
$\phi(a) = 1$, $\phi(X \backslash O) = \{0\}$ and $\varphi = \phi \otimes y \in M$.
Take $(B,\alpha) \in \cB \times I$ such that
\begin{equation}\label{ExpCX-Eq7}
\sup_{n \in \N} \Big(\frac{1}{2n+1} \sum_{j=-n}^n \|(C_{w,f})^j(\varphi)\|_{B,\alpha}\Big) = \infty.
\end{equation}
For each $n \in \Z$, we have that
\begin{equation}\label{ExpCX-Eq8}
\|(C_{w,f})^n(\varphi)\|_{B,\alpha} = \|w^{(n)} \cdot (\varphi \circ f^n)\|_{B \cap f^{-n}(O),\alpha} 
  \leq \|w^{(n)}\|_{B \cap f^{-n}(O)} \|y\|_\alpha.
\end{equation}
By (\ref{ExpCX-Eq7}) and (\ref{ExpCX-Eq8}), we obtain
\[
\sup_{n \in \N} \Big(\frac{1}{2n+1} \sum_{j=-n}^n \|w^{(j)}\|_{B \cap f^{-j}(O)}\Big) = \infty,
\]
as it was to be shown.

Conversely, suppose that the condition in (b) holds.
Given $\varphi \in M \backslash \{0\}$, there exist $C > 0$, $\alpha \in I$ and $O \in \cO$ such that
\[
\|\varphi(x)\|_\alpha \geq C \ \ \text{ for all } x \in O.
\]
By hypothesis, there exists $B \in \cB$ such that
\begin{equation}\label{ExpCX-Eq9}
\sup_{n \in \N} \Big(\frac{1}{2n+1} \sum_{j=-n}^n \|w^{(j)}\|_{B \cap f^{-j}(O)}\Big) = \infty.
\end{equation}
For each $n \in \Z$, we have that
\begin{equation}\label{ExpCX-Eq10}
\|(C_{w,f})^n(\varphi)\|_{B,\alpha} \geq \|w^{(n)} \cdot (\varphi \circ f^n)\|_{B \cap f^{-n}(O),\alpha} 
  \geq C\, \|w^{(n)}\|_{B \cap f^{-n}(O)}.
\end{equation}
By (\ref{ExpCX-Eq9}) and (\ref{ExpCX-Eq10}), we obtain
\[
\sup_{n \in \N} \Big(\frac{1}{2n+1} \sum_{j=-n}^n \|(C_{w,f})^j(\varphi)\|_{B,\alpha}\Big) = \infty,
\]
which shows that $C_{w,f}$ is average expansive. 

\smallskip\noindent
(a): The proof is similar to that of (b).
\end{proof}

\end{document}
